\section{Conclusion}
\label{sec:conclusion}

ALAN is a wind-propelled sailboat built to carry out waypoint missions with
little energy beyond what its sensors and radio require. The team built a
fiberglass displacement hull from a printed mold, a 304 stainless keel with
adjustable ballast canisters, a servo-driven spade rudder, and a
self-trimming wingsail whose tail is driven directly by a small servo. The
firmware divides the STM32H755's two cores by timing requirement: the CM7 runs
sensing, sail planning, PID control, and eleven operating modes under
FreeRTOS, and the CM4 runs the LoRa link. The cores share data through a
versioned, lock-free shared-memory interface. Single-writer arbiters with
timeouts keep each actuator safe if a controller stalls, and a
firmware-in-the-loop simulator built from the CAD let the team develop the
control logic before water time was available. In its first field test the
boat sailed \SI{375}{\meter} under semi-autonomous control and held
reachable headings within 4 to \SI{8}{\degree}, while the self-trimming wing
followed the wind with tail authority to spare. The boat pointed within
\SI{42}{\degree} to \SI{46}{\degree} of the wind, and the radio delivered
95.9\% of packets. The same test confirmed the limits that the remaining
work addresses: without the tack planner, the boat cannot turn through the
wind, and the command acknowledgment scheme loses results during rapid
command sequences.
